\section{Conclusion}

We presented \systemname{}, a systems pattern for turning long-horizon coding-agent work into persistent, auditable optimization loops. The core lesson is simple: for long tasks, roadmaps are weights. They carry the slow project knowledge that local adapters and patches should update against, not overwrite. Across seven CLI-specific skills and one multi-agent orchestrator, \ours{} uses the same non-parametric deep-learning analogy: roadmaps act as slow backbones, active task files act as fast adapters, check phases produce backward signals, local patches guide the next optimize pass, and failure banks preserve reusable error memory. We further described hook, daemon, OS-level cron, and application-level loop boundaries, and included \texttt{moa-loop} as the natural extension from one persistent loop to a mixture-of-agents computation graph.

The immediate value of \ours{} is practical: it gives coding agents a concrete way to continue long-running work without relying on hidden chat context. Future work should add benchmark-scale evaluation, stronger state schemas, richer visual monitoring, and adapters for hosted agent runtimes where file-system persistence is not available by default.
